\documentclass[letterpaper,12pt]{article}
\usepackage[utf8]{inputenc}
\usepackage[spanish]{babel}
\usepackage[T1]{fontenc}
\usepackage{geometry}
\usepackage{graphicx}
\usepackage{enumitem}
\usepackage{color}
\usepackage{tabularx}
\usepackage{float}
\usepackage{caption}
\graphicspath{{img/}}

% Configuración de márgenes
\geometry{
    left=2.5cm,
    right=2.5cm,
    top=2.5cm,
    bottom=2.5cm
}

\begin{document}

%%%%%% ENCABEZADO %%%%%%%

\noindent \textbf{Universidad Central de Venezuela}\\
\textbf{Facultad de Ciencias}\\
\textbf{Escuela de Computación}\\
\textbf{Organización y Estructura del Computador II}\\
Semestre 01-2026

\vspace{0.3cm}
\begin{center}
{\textbf{Práctica 3: Programación en Ensamblador de RISC-V}}
\end{center}
\vspace{0.5cm}

%%%%%% CONTENIDO %%%%%%%

\begin{enumerate}

    %%%%% PREGUNTA 1 %%%%%
    \item Indica el resultado, contenido final de los registros y posiciones de memoria luego de ejecutar las siguientes instrucciones:
    \vspace{0.3cm}
    
    \begin{tabular}{ll@{\hspace{2.5cm}}ll}
        a) & \texttt{add s1, s1, s2} & g) & \texttt{mv s2, s4} \\
        b) & \texttt{addi s3, s2, 2} & h) & \texttt{lw s1, 0(s4)} \\
        c) & \texttt{sub s4, s3, x0} & i) & \texttt{lw s2, 4(s5)} \\
        d) & \texttt{andi s2, s3, -1453} & j) & \texttt{and s5, s1, s3} \\
        e) & \texttt{sll s4, s2, s5} & k) & \texttt{sw s3, 0(s5)} \\
        f) & \texttt{or s1, s1, s2} & l) & \texttt{sw s4, 4(s4)}
    \end{tabular}
    
    \vspace{0.5cm}
    Los registros y memoria tienen los siguientes contenidos en el momento inicial:
    \vspace{0.3cm}

    \begin{table}[H]
        \centering
        \begin{minipage}[t]{0.45\textwidth}
            \centering
            \begin{tabular}{|c|c|}
                \hline
                \multicolumn{2}{|c|}{\textbf{Registros}} \\
                \hline
                \textbf{s1} & \texttt{0x00000016} \\
                \hline
                \textbf{s2} & \texttt{0x00000054} \\
                \hline
                \textbf{s3} & \texttt{0xffffffff} \\
                \hline
                \textbf{s4} & \texttt{0x00000000} \\
                \hline
                \textbf{s5} & \texttt{0x00000004} \\
                \hline
            \end{tabular}
        \end{minipage}
        \hfill
        \begin{minipage}[t]{0.45\textwidth}
            \centering
            \begin{tabular}{|c|c|}
                \hline
                \multicolumn{2}{|c|}{\textbf{Memoria}} \\
                \hline
                \textbf{0x00} & \texttt{0x03393826} \\
                \hline
                \textbf{0x04} & \texttt{0xea0063af} \\
                \hline
                \textbf{0x08} & \texttt{0x17fa8912} \\
                \hline
                \textbf{0x0c} & \texttt{0xbc983304} \\
                \hline
                \textbf{0x10} & \texttt{0x534b4aaa} \\
                \hline
                \textbf{0x14} & \texttt{0x534b4aaa} \\
                \hline
            \end{tabular}
        \end{minipage}
    \end{table}

    \vspace{0.5cm}

    %%%%% PREGUNTA 2 %%%%%
    \item Escribe un programa en lenguaje ensamblador que implemente el siguiente bloque \texttt{IF-THEN-ELSE}. Usa la sección \texttt{.data} para asignar algún valor inicial a las variables.
    \begin{verbatim}
    if (x >= y) {
        x = x + 2; 
        y = y + 2;
    } else {
        x = x - 2;
        y = y - 2;
    }
    \end{verbatim}

    \pagebreak

    %%%%% PREGUNTA 3 %%%%%
    \item Escribe un programa en lenguaje ensamblador que implemente el siguiente código. Usa la sección \texttt{.bss} para reservar espacio en memoria para las variables.
    \begin{verbatim}
    n = 5;
    fant = 1;
    f = 1;
    i = 2;
    while (i <= n) {
        faux = f;
        f = f + fant; 
        fant = faux; 
        i = i + 1;
    }
    \end{verbatim}

    \vspace{0.5cm}

    %%%%% PREGUNTA 4 %%%%%
    \item Cada una de las siguientes instrucciones del repertorio de ensamblador de RISC-V, ¿con qué secuencia de 32 bits se corresponde en código máquina?
    \begin{enumerate}[label=\alph*)]
        \item \texttt{lw t0, 0(t2)}
        \item \texttt{bge t1, t0, 0x2C}
        \item \texttt{sw t1, 0(t4)}
    \end{enumerate}

    \vspace{0.5cm}

    %%%%% PREGUNTA 5 %%%%%
    \item Cada una de las siguientes secuencias de 32 bits, ¿con qué instrucciones del repertorio de ensamblador de RISC-V se corresponden?
    \begin{enumerate}[label=\alph*)]
        \item \texttt{0x03528B33}
        \item \texttt{0x00190913}
        \item \texttt{0x0000006F}
    \end{enumerate}

\end{enumerate}

\end{document}
